%! Sinusoidal Functions — Unit 3, Lesson 3.5 — Exit Ticket (Answer Key)
\documentclass[10pt]{article}
\usepackage{precalculus-article}
\usepackage{precalculus-key}
\newcommand{\TallMath}[1]{$\displaystyle #1\rule[-1.4em]{0pt}{3.2em}$}

\begin{document}

\pageheader{Unit 3, Lesson 3.5}{Exit Ticket --- Answer Key}
\namedateperiod

\vspace{0.18in}

\begin{notesbox}{Exit Ticket: Sinusoidal Functions --- Key}

\textbf{1.} A sinusoidal function has a maximum value of $7$ and a minimum value
of $1$.

\medskip
(a) What is the amplitude of this function?

\vspace{0.08in}
\begin{tabular}{@{}l@{}}
  Amplitude $= \dfrac{\ans{7} - \ans{1}}{2} = $ \ans{3}
\end{tabular}

\vspace{0.18in}

(b) What is the midline of this function?  Write your answer as an equation.

\vspace{0.08in}
\begin{tabular}{@{}l@{}}
  Midline: $y = \dfrac{\ans{7} + \ans{1}}{2} = $ \ans{$y = 4$}
\end{tabular}

\vspace{0.25in}

\textbf{2.} The period of a sinusoidal function is $\dfrac{\pi}{2}$.
What is the frequency of this function?

\vspace{0.1in}
\begin{tabular}{@{}l@{}}
  Frequency $= \dfrac{1}{\ans{$\pi/2$}} = $ \ans{$\dfrac{2}{\pi}$}
\end{tabular}

\end{notesbox}

\end{document}
